% ==========================================================================
%  NOTE: Working notes for the thesis, not the final chapter. Documents the
%  proof of concept in whole_apps/honeybee_poc/: applying statesnap to the REAL,
%  unflattened BioDT honeybee app, with a Plumber backend, communicating via
%  statesnap JSON over HTTP. Verify claims against the committed code before
%  citing. Session of 2026-06-12.
% ==========================================================================

\documentclass[12pt,a4paper]{article}
\usepackage[utf8]{inputenc}
\usepackage[T1]{fontenc}
\usepackage{geometry}
\geometry{margin=2.5cm}
\usepackage{hyperref}
\usepackage{booktabs}
\usepackage{enumitem}
\usepackage{listings}
\usepackage{xcolor}
\lstset{basicstyle=\ttfamily\footnotesize,breaklines=true,frame=single,
        columns=fullflexible,keepspaces=true}

\title{\textbf{Applying \texttt{statesnap} to a Modular Scientific Shiny App\\
without Flattening It}\\[0.4em]
\large Proof of concept: BioDT Honeybee frontend $\leftrightarrow$ Plumber backend
over \texttt{statesnap} JSON}
\author{Iva Gunjaca \\ Master of Software Engineering, University of Amsterdam}
\date{\today}

\begin{document}
\maketitle

\section*{Attribution}
The honeybee (Beekeeper pDT) application is from the \textbf{LTER LIFE} project
(BioDT, \texttt{github.com/BioDT/biodt-shiny}) and is used here with their
permission for this thesis. All credit for the original application and its
scientific model belongs to the LTER LIFE / BioDT authors. The compute backend
in this PoC is a \emph{simulated} stand-in for \texttt{ghcr.io/biodt/beehave},
not the real NetLogo model.

% --------------------------------------------------------------------------
\section{The problem this PoC solves}
% --------------------------------------------------------------------------

The feasibility note (\texttt{biodt\_feasibility.tex}) identified a blocking gap:
\texttt{statesnap}'s \texttt{capture\_state()} reads \texttt{names(input)} at the
top level of a Shiny session, so it cannot see widgets that live inside Shiny
\textbf{module namespaces} (\texttt{ns("x")}). The earlier honeybee integration
(\texttt{ai\_integration\_experiment.tex}) sidestepped this by \emph{flattening}
the app into a single plain-Shiny file. That works for a benchmark but throws
away exactly what makes the real BioDT app interesting: it is a deeply modular
Rhino/\texttt{box} application.

\textbf{This PoC keeps the real app modular and adapts \texttt{statesnap} to it.}
The question it answers: \emph{can state be captured and restored on a
namespaced, modular scientific app without rewriting the app?}

% --------------------------------------------------------------------------
\section{Key idea: capture the MODEL, not the inputs}
% --------------------------------------------------------------------------

The insight is that the BioDT honeybee app already separates \emph{model} from
\emph{widgets}. Its parent module (\texttt{honeybee\_beekeeper.R}) composes the
child modules into a small set of \textbf{namespace-free reactive model
fragments}, which the children \emph{return} from their servers:

\begin{itemize}[nosep]
  \item \texttt{coordinates()} \quad from \texttt{beekeeper\_map}
  \item \texttt{parameters()} \quad from \texttt{beekeeper\_param}
        (a list of two data frames: \texttt{parameters}, \texttt{simulation})
  \item \texttt{lookup()} \quad from \texttt{beekeeper\_lookup}
  \item \texttt{experiment\_list()} \quad results, from \texttt{beekeeper\_runsimulation}
\end{itemize}

Each child owns its \texttt{ns()} internally and hands back a plain value. So
instead of asking \texttt{statesnap} to introspect namespaced inputs, we capture
the \textbf{composed model}:

\begin{lstlisting}[language=R]
# input = list()  -> capture_state never reads names(input),
# so module namespaces are irrelevant. The model is the single
# registered, JSON-serialisable value.
json <- capture_state(input = list(), model = model)
\end{lstlisting}

This is the whole trick: by passing \texttt{input = list()} and registering the
domain model as a named value, \texttt{statesnap} operates at the
\textbf{model layer}. No flattening, no change to how the child modules use
namespaces, no change to \texttt{statesnap} itself.

% --------------------------------------------------------------------------
\section{Architecture: statesnap as both checkpoint and wire format}
% --------------------------------------------------------------------------

The same \texttt{statesnap} JSON serves two purposes at once:

\begin{enumerate}[nosep]
  \item \textbf{Transport.} The frontend POSTs the model checkpoint to a Plumber
    backend over HTTP; the backend returns the result as a checkpoint. No
    Spring, no Kafka, no Redis -- two R processes exchanging \texttt{statesnap}
    JSON.
  \item \textbf{Reproducibility.} The same JSON, saved to a file, reopens a run
    later -- including the \emph{computed, non-deterministic} result.
\end{enumerate}

\begin{center}
\footnotesize
\begin{tabular}{@{}lcl@{}}
\toprule
\textbf{Frontend (real BioDT app)} & & \textbf{Backend (Plumber)} \\
\midrule
child modules return model fragments & & \\
$\downarrow$ \texttt{build\_model()} & & \\
$\downarrow$ \texttt{capture\_state(input=list(), model)} & & \\
\texttt{POST /run} (JSON) & $\longrightarrow$ & \texttt{restore\_state(model=...)} \\
 & & $\downarrow$ \texttt{run\_beehave\_sim()} (simulated) \\
\texttt{result\_from\_json()} & $\longleftarrow$ & \texttt{capture\_state(result=...)} \\
$\downarrow$ write BioDT result CSV & & \\
unchanged \texttt{beekeeper\_plot} renders & & \\
\bottomrule
\end{tabular}
\end{center}

% --------------------------------------------------------------------------
\section{What was changed in the real app (minimal footprint)}
% --------------------------------------------------------------------------

Only \textbf{five} files of the upstream app are touched -- two new, three
patched. Every other module (map, param, lookup, plot) runs unchanged.

\begin{table}[h]\centering\footnotesize
\begin{tabular}{@{}llp{6.2cm}@{}}
\toprule
\textbf{File} & \textbf{Change} & \textbf{Purpose} \\
\midrule
\texttt{logic/honeybee/statesnap\_adapter.R} & new &
  model$\leftrightarrow$JSON$\leftrightarrow$HTTP helpers; the namespace-proof
  capture/restore. \\
\texttt{view/honeybee/beekeeper\_checkpoint.R} & new &
  save/restore-run UI card at the model level. \\
\texttt{view/honeybee/beekeeper\_runsimulation.R} & patched &
  adds a \texttt{"plumber"} executor branch alongside docker/k8s. \\
\texttt{view/honeybee/honeybee\_beekeeper.R} & patched &
  wires the checkpoint module into the parent. \\
\texttt{config.yml} & patched &
  adds a \texttt{poc} profile (\texttt{executor: plumber}). \\
\bottomrule
\end{tabular}
\end{table}

The new executor branch composes the model from the existing reactives and ships
it -- it does not reach into any module's internals:

\begin{lstlisting}[language=R]
} else if (config$get("executor") == "plumber") {
  model <- build_model(coordinates(), parameters(), lookup())
  reply <- run_on_backend(model, base_url = config$get("backend_url"))
  result_to_biodt_csv(reply$result, run_dir)  # unchanged plot reads this
}
\end{lstlisting}

% --------------------------------------------------------------------------
\section{A real bug found during verification (worth reporting)}
% --------------------------------------------------------------------------

The first roundtrip ``worked'' but silently used \emph{default} parameters: the
restored model had empty fields. Cause: \texttt{jsonlite} serialises a
\texttt{data.frame} \textbf{row-oriented}, so after \texttt{restore} the frames
came back as lists of row-objects and column access (\texttt{model\$coordinates\$lat})
returned \texttt{NULL}. The simulation then fell back to defaults -- a dangerous
\emph{silent} failure, not a crash.

\textbf{Fix} (in \texttt{build\_model}): coerce each data frame to a
\textbf{column-oriented} list with \texttt{as.list()} before capture, so columns
round-trip as named vectors. After the fix, the model survives intact
(\texttt{lat}=52.37, \texttt{N\_INITIAL\_BEES}=12000, lookup nectar preserved)
and the simulation reacts to the actual inputs.

\noindent\emph{Thesis relevance:} this is a concrete, citable instance of the
serialisation subtleties that arise when state crosses a process boundary -- and
a reminder that ``it ran'' is not ``it was correct.'' It surfaced only because
verification checked the \emph{restored values}, not just that a run completed.

\noindent\textbf{Update (2026-06-15): fix promoted into the library.} The
\texttt{as.list()} coercion described here was an \emph{application-level}
workaround inside \texttt{build\_model}. Because the same row-oriented hazard
affects \emph{any} use of \texttt{statesnap} that round-trips a data frame, the
fix has since been moved into the library itself: \texttt{capture\_state()} now
detects data frames at any depth and stores them column-oriented, and
\texttt{restore\_state()} rebuilds them with atomic column types, \texttt{NA}s and
character row names preserved (non-atomic classes such as \texttt{factor} or
\texttt{Date} degrade to their JSON scalar form --- the documented boundary). The PoC's manual coercion is consequently no longer
required --- a plain \texttt{capture\_state()} on the model would now preserve the
tables --- though it remains harmless. See the \texttt{statesnap} library write-up
(``Module-namespace and data-frame support'') for the general implementation.

% --------------------------------------------------------------------------
\section{Verification}
% --------------------------------------------------------------------------

\begin{itemize}[nosep]
  \item All six R artifacts \texttt{parse()} cleanly (R 4.3.3).
  \item \texttt{statesnap} installed; full model$\to$JSON$\to$restore$\to$simulate
    $\to$result-JSON$\to$extract roundtrip executed end-to-end (without HTTP).
  \item The captured model JSON contains \emph{no} namespaced inputs
    (\texttt{"inputs":\{\}}) -- confirming model-layer capture.
  \item Restored result carries BioDT's exact plot columns
    (\texttt{date}, \texttt{weather}, \texttt{TotalIHbees + TotalForagers},
    \texttt{(honeyEnergyStore / ( ENERGY\_HONEY\_per\_g * 1000 ))}).
  \item Output is non-deterministic across seeds -- the property that makes
    capturing the \emph{computed} result (not inputs) necessary for
    reproducibility.
  \item \textbf{Not verified:} the full app was not launched under Docker in this
    environment (the BioDT base image and raster data are heavy); the Shiny UI
    integration is verified by parse + the model-contract roundtrip, and is
    intended to be run dockerised on the author's laptop via the provided
    \texttt{docker-compose.yml}.
\end{itemize}

% --------------------------------------------------------------------------
\section{Why this matters for the thesis argument}
% --------------------------------------------------------------------------

\begin{itemize}[nosep]
  \item It closes the namespace gap the feasibility note flagged, \emph{without}
    flattening and \emph{without} modifying \texttt{statesnap} -- the seam is the
    app's own model/view separation. (The library has since \emph{also} gained
    native namespace and data-frame support, so this gap is now closed at two
    levels: the model-layer pattern shown here, and a direct namespaced capture
    --- see the updates later in this document.)
  \item It shows \texttt{statesnap} JSON doubling as a transport between a Shiny
    frontend and a compute backend, i.e. a lightweight alternative to the
    Spring/Kafka relay for the single-app case.
  \item The five-file, mostly-additive footprint is evidence that a
    well-structured modular pDT can gain reproducible, shareable state cheaply.
  \item The serialisation bug is an honest, instructive negative result about
    crossing process boundaries with R objects.
\end{itemize}

\paragraph{Generalisation.} The model-layer technique requires only that a
modular app exposes its model as return values from module servers (as BioDT
does). At the time of writing the PoC, an app that buried its state in namespaced
inputs with \emph{no} model layer would have needed either a namespace-aware
\texttt{capture\_state} or a small refactor to expose a model.

\noindent\textbf{Update (2026-06-15):} that namespace-aware path now exists.
\texttt{capture\_state()} / \texttt{restore\_state()} take an optional \texttt{ns}
argument that slices a single module's inputs out by namespace (storing them
under local names and re-qualifying on restore), and capture/restore invoked
inside a module's own server works unmodified because Shiny scopes the module
\texttt{input}. So the model-layer approach shown here and a direct namespaced
capture are now \emph{two supported options} rather than the model layer being the
only one. The model-layer route remains attractive when an app already exposes a
clean model, because it keeps the checkpoint independent of widget identity.

\paragraph{Next steps.} Run the stack under Docker end-to-end and screenshot the
save/restore of a non-deterministic run; and swap the simulated backend for the
real \texttt{ghcr.io/biodt/beehave} container behind the same \texttt{statesnap}
contract. (The previously-listed ``extend \texttt{capture\_state} with namespace
awareness'' item is now done --- see the update above.)

\end{document}
